\section{Verification of the Numerical Hypotheses}
\label{app:verify}

This appendix proves Theorem~\ref{thm:N-holds}.

Throughout, $a_0,d,\xiL,\xia,\xib,\xiR$, $H$, $\chi$, $\prof$, $\Omega_G$,
$\psi_\ell$, $\Evans_\ell$, $\Thl{\ell}$ and $\epsst$ are as in
Definition~I.\ref{I-def:ring} and Hypothesis~I.\ref{I-hyp:N}, and $G\ge0$. Only the
properties of $\chi$ listed in Definition~I.\ref{I-def:ring} are used. We identify $\C$ with $\R^2$, write $|Y|_w=\max(\ell|u|,|v|)$ for
$\ell\ge1$ and $Y=(u,v)\in\C^2$, and put
\[
x_a=\frac{\xia-a_0}{d},\quad y_b=-\frac{\xib}{d},\quad \Gmax=\frac1{50},\quad \Gmax'=d,\quad
\Rect=\Big[\frac9{200},3\Big]\times[-5,3].
\]
Definition~I.\ref{I-def:ring} gives exactly $\xia-a_0=-\frac12+2^{-54}$ and
$d=\frac1{50}(1+6\cdot2^{-58})$. So $d>\frac1{50}$, whence the range
$[0,\Gmax']$ of the computations contains the range $[0,\Gmax]$ of $G$ in
\textup{(N1)} and \textup{(N2)}, and $50(1-6\cdot2^{-58})<1/d<50$, so
$x_a=-(\frac12-2^{-54})/d$ and $y_b=-1/(2d)$ satisfy
\begin{equation}
\label{eq:xayb}
-25<x_a<-25+4\cdot10^{-15},\qquad -25<y_b<-25+10^{-15}.
\end{equation}
For a continuous closed curve $c\colon[0,1]\to\C\setminus\{0\}$,
$\wind(c,0)=(\theta(1)-\theta(0))/2\pi$, where $\theta$ is any continuous
function with $c=|c|e^{i\theta}$. If $\Gamma$ is the boundary of a closed
rectangle in $\{\Re\lambda>0\}$ and $\Evans_3(\cdot;G)\ne0$ on $\Gamma$,
$\wind(\Evans_3(\Gamma;G),0)$ is the winding number of
$t\mapsto\Evans_3(\gamma(t);G)$, where $\gamma$ traverses $\Gamma$ once
counterclockwise.

\subsection{Ball Arithmetic and the Trusted Base}
\label{app:arith}

The computations use \textbf{ball arithmetic}: a real number is enclosed in a
ball $[m\pm r]=[m-r,m+r]$ whose midpoint $m$ and radius $r\ge0$ are binary
floating-point numbers, and a complex number in a box $X+iY$, the product of two
real balls. We use the implementation of Johansson \cite{Johansson2017}, part of
FLINT \cite{FLINT} from version~3 on, in FLINT~3.6.0 as bundled with
version~0.9.0 of its Python interface python-flint \cite{pythonflint}, at a
working precision of $128$ bits. For a real ball $X$, $\mathrm{lo}(X)$ and
$\mathrm{up}(X)$ are its endpoints, which are dyadic rationals; for a box $Z$,
$|Z|^{\rm up}$ is a dyadic upper bound of $|z|$ over $z\in Z$; for a ball $\rho$,
$\mathrm{box}(\rho)=[-\mathrm{up}(\rho),\mathrm{up}(\rho)]
+i[-\mathrm{up}(\rho),\mathrm{up}(\rho)]$, which contains the closed disk of
radius $\rho'$ about $0$ for every $\rho'\in\rho$ with $\rho'\ge0$;
$\mathrm{hull}(X,X')$ is a ball computed from $\mathrm{lo}(X)$ and
$\mathrm{up}(X')$; and $X\cup X'$ is the computed union of two real balls, again
a ball. Every enclosure rests on the following \textbf{containment property}.
\begin{enumerate}[label=\textup{(CP)}, leftmargin=*]
\item For every arithmetic operation $f$ that we use (sums, differences, products
  and quotients of real balls and of boxes, integer powers of a real ball, the
  maximum of two real balls, and the exponential and the square root of a real
  ball), and for all input balls $X_1,\dots,X_k$, the output ball contains
  $f(x_1,\dots,x_k)$ for every choice of $x_j\in X_j$ at which $f$ is defined; a
  ball with infinite radius, or with an undefined midpoint, is regarded as the
  whole real line. A ball built from an integer contains that integer, a ball
  built from an exact midpoint $m$ and an exact radius $r$ is $[m-r,m+r]$, and a
  box built from two real balls is their product. Moreover
  $\mathrm{lo}(X)\le x\le\mathrm{up}(X)$ for every $x\in X$, $|z|\le|Z|^{\rm up}$ for
  every $z\in Z$, $\mathrm{hull}(X,X')$ contains every real number between
  $\mathrm{lo}(X)$ and $\mathrm{up}(X')$, $X\cup X'$ contains $X$, $X'$ and every
  real number between a point of $X$ and a point of $X'$, the computed
  intersection of a ball $X$ with $[0,\infty)$ contains $X\cap[0,\infty)$, and a
  comparison such as $X>0$ or $X\le c$ returns true only if every point of $X$
  satisfies it.
\end{enumerate}
The data of Definition~I.\ref{I-def:ring}, the contour points, and the endpoints of
the intervals of $G$ and the rectangles of Section~\ref{app:runs} are rational.
They are held exactly, as pairs of integers, in the arbitrary-precision integer
arithmetic of CPython, and every comparison between them is made there. A
rational $n/k$ in lowest terms, $k\ge1$, enters ball arithmetic as follows. If
$k=2^e$ and $|n|<2^{125}$, it enters as the product of the ball built from $n$ and
the power $-e$ of the ball built from $2$; the program checks that this product
has radius $0$, so by (CP) it is the point $n/k$. Otherwise it enters as the
quotient of the balls built from $n$ and $k$, which contains $n/k$ by (CP). No
hardware floating-point number of the programming language enters any enclosure:
every ball is built from integers or from other balls.

The proof of Proposition~\ref{prop:runs} is a computation. It is correct provided
the following \textbf{trusted base} holds.
\begin{enumerate}[label=\textup{(TB\arabic*)}, leftmargin=*]
\item The ball arithmetic of FLINT~3.6.0, as bundled with python-flint~0.9.0,
  satisfies (CP) for the operations listed there.
\item CPython integer arithmetic, and the rational arithmetic of its standard
  module \texttt{fractions}, are correct.
\item The program runs as written, with its assertions enabled, and its routines
  carry out the enclosure procedure of Section~\ref{app:procedure} and the
  certifications and runs of Section~\ref{app:runs}.
\end{enumerate}
We did not inspect FLINT's internal error analysis; if FLINT uses hardware
floating-point arithmetic internally, its correctness is part of (TB1). The
function \texttt{evans} of the module \texttt{ringcert.evans} follows steps
(i)--(iv) of Section~\ref{app:procedure} line by line, and the part of (TB3) that
concerns it is established by reading it against them; the tests of the package
detect some departures from the procedure (dropping the Taylor remainder in
step~(iii)), but not all (omitting the a priori enlargement that defines
$P^{\rm st}$ and $P'^{\rm st}$ in step~(iii)), and do not replace this reading. A certification (Section~\ref{app:runs})
is carried out by the function \texttt{certify\_cell} of the module
\texttt{ringcert.winding}, which follows the description of the certification in
Section~\ref{app:runs} line by line, and a run
by the function \texttt{certify\_cell\_adaptive} of the same module. Everything
else in this paper has a written proof that uses no machine computation.

\subsection{Two Elementary Lemmas}

\begin{lemma}
\label{lem:end-regions}
The following two estimates hold.
\begin{enumerate}[label=\textup{(\arabic*)}]
\item We have $\int_{-\infty}^{\xia}|\prof'|\dd\xi\le\frac94e^{2x_a}$ and
  $\int_{\xib}^{\infty}|\prof'|\dd\xi\le\frac94e^{2y_b}$, and both bounds are less
  than $3e^{-50}$.
\item Let $\ell\ge1$, $G\ge0$ and $\Re\lambda>0$, let $\psi=\psi_\ell(\cdot;\lambda,G)$,
  and let $[\alpha,\beta]$ be one of the intervals $[\xiL,\xia]$ and $[\xib,\xiR]$.
  Let $\psi_0$ solve $\psi_0''=\ell^2\psi_0$ with $\psi_0(\alpha)=\psi(\alpha)$ and
  $\psi_0'(\alpha)=\psi'(\alpha)$, and put $Y=(\psi,\psi')$, $Y_0=(\psi_0,\psi_0')$
  and $\epsend=3e^{-50}/\Re\lambda$. Then, for every $\xi\in[\alpha,\beta]$,
  \[
  |Y(\xi)-Y_0(\xi)|_w\le e^{\ell(\xi-\alpha)}\,|Y(\alpha)|_w\,\big(e^{\epsend}-1\big).
  \]
\end{enumerate}
\end{lemma}
\begin{proof}
(1) Write $A(\xi)=H(\frac{\xi-a_0}d)$, $B(\xi)=H(-\frac\xi d)$, so that
$\prof=\chi AB$, and $x=\frac{\xi-a_0}d$, $y=-\frac\xi d$. Since
$H'=\frac12\operatorname{sech}^2$, $A'=\frac1{2d}\operatorname{sech}^2x\ge0$ and
$B'=-\frac1{2d}\operatorname{sech}^2y\le0$; we use $H(x)\le e^{2x}$,
$\operatorname{sech}^2x\le4e^{2x}$, $0\le H\le1$ and $\operatorname{sech}^2\le1$. On
$(-\infty,\xia]$, $\prof'=\chi'AB+\chi(A'B+AB')$, $\chi$ is nondecreasing from
$\chi(\xiL)=0$ to $\chi(\xia)=1$, and $AB\le A\le A(\xia)=H(x_a)$ since $A$ is
increasing; so $\int_{-\infty}^{\xia}|\chi'|AB\le H(x_a)\le e^{2x_a}$. Next,
$0\le\chi\le1$, $|A'B+AB'|\le\frac1{2d}(\operatorname{sech}^2x+H(x)\operatorname{sech}^2y)
\le\frac5{2d}e^{2x}$ and $\int_{-\infty}^{\xia}e^{2(\xi-a_0)/d}\dd\xi=\frac d2e^{2x_a}$;
adding gives $\frac94e^{2x_a}$. On $[\xib,\infty)$ exchange the roles of $x$ and
$y$: $\chi$ is nonincreasing from $1$ to $0$, $AB\le B\le B(\xib)=H(y_b)\le e^{2y_b}$,
$|A'B+AB'|\le\frac1{2d}(\operatorname{sech}^2x\,H(y)+\operatorname{sech}^2y)\le\frac5{2d}e^{2y}$
and $\int_{\xib}^\infty e^{-2\xi/d}\dd\xi=\frac d2e^{2y_b}$. By \eqref{eq:xayb}, both
bounds are at most $\frac94e^{-50}e^{8\cdot10^{-15}}<3e^{-50}$, as
$e^{8\cdot10^{-15}}<\frac43$.

(2) The vector $Z=(\ell\psi,\psi')$ satisfies $Z'=A_0Z+\mathcal QZ$ with
$A_0=\begin{psmallmatrix}0&\ell\\ \ell&0\end{psmallmatrix}$ and
$\mathcal Q=\begin{psmallmatrix}0&0\\ \delta/\ell&0\end{psmallmatrix}$, where
$\delta=i\ell\prof'/(\lambda+i\ell\Omega_G)$, and $Z_0=(\ell\psi_0,\psi_0')$ satisfies
$Z_0'=A_0Z_0$; and $|Y|_w=\|Z\|_\infty$, the maximum of the moduli of the
components. Since $\Omega_G$ is real, $|\lambda+i\ell\Omega_G|\ge\Re\lambda$, so
the $\|\cdot\|_\infty$ operator norm of $\mathcal Q$ is
$b=|\prof'|/|\lambda+i\ell\Omega_G|\le|\prof'|/\Re\lambda$, and
$\int_\alpha^\beta b\le\epsend$ by (1). For $t\ge0$,
$e^{tA_0}=\begin{psmallmatrix}\cosh\ell t&\sinh\ell t\\ \sinh\ell t&\cosh\ell t\end{psmallmatrix}$
has norm $e^{\ell t}$. By variation of constants,
$Z(\xi)-Z_0(\xi)=\int_\alpha^\xi e^{(\xi-\eta)A_0}\mathcal Q(\eta)Z(\eta)\dd\eta$. With
$u(\xi)=e^{-\ell(\xi-\alpha)}\|Z(\xi)-Z_0(\xi)\|_\infty$, $c=\|Z(\alpha)\|_\infty$ and
$\|Z_0(\eta)\|_\infty\le e^{\ell(\eta-\alpha)}c$, this gives
$u(\xi)\le\int_\alpha^\xi b(\eta)(c+u(\eta))\dd\eta$, so by Gronwall's inequality
$c+u(\xi)\le c\exp(\int_\alpha^\xi b)$ and $u(\xi)\le c\,(e^{\epsend}-1)$.
\end{proof}

\begin{lemma}
\label{lem:apriori}
Let $0\le G\le\Gmax'$. Every zero $\lambda$ of $\Evans_3(\cdot;G)$ with
$\Re\lambda\ge\frac9{200}$ lies in $\Rect$.
\end{lemma}
\begin{proof}
Let $\ell=3$, let $\lambda$ be a zero of $\Evans_\ell(\cdot;G)$ with
$\Re\lambda\ge\frac9{200}$, so $\Re\lambda>0$, and put $TV=\int_{\R}|\prof'|\dd\xi$ and
$\Omega_{\max}=\frac12+\Gmax'e^{-2\xiL}$. On $[\xia,\xib]$, $\chi=1$, so, with $A,B$ as in the proof of
Lemma~\ref{lem:end-regions}, $\prof'=A'B+AB'$ with $A'\ge0\ge B'$ and $0\le A,B\le1$.
Hence $\int_{\xia}^{\xib}|\prof'|\le\int_{\xia}^{\xib}(A'-B')
=A(\xib)-A(\xia)+B(\xia)-B(\xib)<2$, and Lemma~\ref{lem:end-regions}(1) gives
$TV<2+6e^{-50}$. Since $e>2.7$ and $2.7^7>1000$, $e^{50}>2.7^{49}=(2.7^7)^7>10^{21}$, so
$6e^{-50}<10^{-20}$; and
$\log2=\sum_{k\ge1}\frac1{k2^k}>\sum_{k=1}^6\frac1{k2^k}>0.69$. So
$TV/\log2<(2+10^{-20})/0.69<2.9$.

For $\xi\in\R$, $0\le\prof\le1$ gives
$0\le\int_{-\infty}^\xi\prof(\xi')e^{2\xi'}\dd\xi'\le\frac12e^{2\xi}$, so
$0\le\Omega_G(\xi)\le\frac12+Ge^{-2\xi}\le\Omega_{\max}$ for $\xi\ge\xiL$. By
Definition~I.\ref{I-def:ring}, $-2\xiL=4528212286209965\cdot2^{-51}<2.011$; $e<2.7183$ gives $e^2<7.3892$, and
$e^x\le1+x+x^2$ ($0\le x\le1$) gives $e^{0.011}<1.01113$, so
$e^{-2\xiL}<7.3892\cdot1.01113<7.4715$. As $d<0.0200001$,
$\Omega_{\max}<\frac12+0.0200001\cdot7.4715<0.6495$.

Let $\psi=\psi_\ell(\cdot;\lambda,G)$ and
$\mu=\min_{\xi\in[\xiL,\xiR]}|\lambda+i\ell\Omega_G(\xi)|$, positive as $\Re\lambda>0$. The proof of
Lemma~\ref{lem:end-regions}(2), run on $[\xiL,\xiR]$ with
$\psi_0(\xi)=e^{\ell(\xi-\xiL)}$ (which has the same data at $\xiL$),
$\int_{\xiL}^{\xiR}b\le TV/\mu=:\varepsilon'$ and $|Y(\xiL)|_w=\ell$, gives
$|Y(\xiR)-Y_0(\xiR)|_w\le e^{\ell(\xiR-\xiL)}\ell\,(e^{\varepsilon'}-1)$. Since
$\psi_0'+\ell\psi_0=2\ell e^{\ell(\xi-\xiL)}$ and
$|(\psi'-\psi_0')+\ell(\psi-\psi_0)|\le2|Y-Y_0|_w$, we get
$|\Evans_\ell(\lambda;G)-2\ell|\le2\ell(e^{\varepsilon'}-1)$. If $\mu>TV/\log2$, then
$\varepsilon'<\log2$, the right side is less than $2\ell$, and
$\Evans_\ell(\lambda;G)\ne0$. As $\lambda$ is a zero, $\mu\le TV/\log2$, and $\mu$ is attained at some
$\xi\in[\xiL,\xiR]$, so $|\lambda+i\ell\,\Omega_G(\xi)|\le TV/\log2$. At that $\xi$,
$\Re\lambda\le TV/\log2<2.9$ and $|\Im\lambda+\ell\Omega_G(\xi)|\le TV/\log2$, so
$0\le\Omega_G(\xi)\le\Omega_{\max}$ gives $-\ell\,\Omega_{\max}-2.9<\Im\lambda<2.9$. As
$\ell=3$, $-3\Omega_{\max}-2.9>-3\cdot0.6495-2.9>-4.85$. Hence
$\lambda\in[\frac9{200},2.9)\times(-4.85,2.9)\subset\Rect$.
\end{proof}

\subsection{The Enclosure Procedure}
\label{app:procedure}

The procedure encloses $\Evans_\ell(\lambda;G)$ for all $\lambda$ in a box and
all $G$ in an interval. It crosses the end regions by
Lemma~\ref{lem:end-regions}, and $[\xia,\xib]$, where $\chi=1$, by a Taylor
method with coefficients and remainders enclosed in ball arithmetic. Fix the
Taylor order $\NT=16$, written $N$ in this subsection, and the points $s_0<s_1<\dots<s_n$, $n=360$, given by
$s_0=\xia$, $s_n=\xib$ and $s_j=\RN(\xia+\frac jn(\xib-\xia))$ for $0<j<n$, where
$\RN(q)$ is the binary64 number nearest to $q$ (ties to even). Only three facts are
used: the $s_j$ are dyadic rationals, they are strictly increasing, and the steps
$h_j=s_{j+1}-s_j$ are exact dyadic rationals.

The input is an integer $\ell\ge1$, a box $\Lambda^{\rm b}=\mathsf X+i\mathsf Y$ and a
real ball $\mathcal G^{\rm b}$. Each operation below is a ball operation; the
procedure \textbf{fails} if a comparison that it requires is not certified as
true, or if its output is not a box of finite radius.
\begin{enumerate}[label=\textup{(\roman*)}]
\item \emph{Coefficient data, independent of $\lambda$ and $G$.} For each $j<n$,
  and for $S=\{s_j\}$ (the \emph{point data}) and for $S=[s_j,s_{j+1}]$ (the
  \emph{step data}), compute
  \[
  \begin{gathered}
  \mathsf A_0=\mathrm{hull}\big(H(\tfrac{\min S-a_0}d),H(\tfrac{\max S-a_0}d)\big),\quad
  \bar{\mathsf A}_0=\mathrm{hull}\big(H(\tfrac{a_0-\max S}d),H(\tfrac{a_0-\min S}d)\big),\\
  \mathsf B_0=\mathrm{hull}\big(H(\tfrac{-\max S}d),H(\tfrac{-\min S}d)\big),\quad
  \bar{\mathsf B}_0=\mathrm{hull}\big(H(\tfrac{\min S}d),H(\tfrac{\max S}d)\big),
  \end{gathered}
  \]
  with $H(x)$ evaluated as $1/(1+\exp(-2x))$; then, for $0\le k\le N$,
  \[
  \begin{gathered}
  \mathsf A_{k+1}=\frac{2}{d(k+1)}\sum_{r=0}^k\mathsf A_r\bar{\mathsf A}_{k-r},\qquad
  \bar{\mathsf A}_{k+1}=-\mathsf A_{k+1},\\
  \mathsf B_{k+1}=-\frac{2}{d(k+1)}\sum_{r=0}^k\mathsf B_r\bar{\mathsf B}_{k-r},\qquad
  \bar{\mathsf B}_{k+1}=-\mathsf B_{k+1},
  \end{gathered}
  \]
  and $\mathsf w_k=\sum_{r\le k}\mathsf A_r\mathsf B_{k-r}$ for $k\le N+1$,
  $\mathsf f_k=(k+1)\mathsf w_{k+1}$, $\mathsf E^\pm_k=(\pm2)^ke^{\pm2S}/k!$ (with $e^{\pm2S}$ a
  ball containing $e^{\pm2\xi}$ for $\xi\in S$) and
  $\mathsf q_k=\sum_{r\le k}\mathsf w_r\mathsf E^+_{k-r}$ for $k\le N$. Then compute the
  balls
  \[
  \gamma_0=\mathrm{hull}\Big(0,\frac{e^{2x_a+2\xia}}{2/d+2}\Big),\qquad
  \gamma_{j+1}=\gamma_j+\Big(\sum_{k<N}\mathsf q_k^{\rm pt}\frac{h_j^{k+1}}{k+1}
  +\mathsf q_N^{\rm st}\frac{h_j^{N+1}}{N+1}\Big)\cap[0,\infty),
  \]
  where the superscripts ${\rm pt}$ and ${\rm st}$ refer to the point and the step
  data of step $j$, and, with $\mathsf g_0=\gamma_j$ for the point data,
  $\mathsf g_0=\mathrm{hull}(\gamma_j,\gamma_{j+1})$ for the step data, and
  $\mathsf g_r=\mathsf q_{r-1}/r$ for $1\le r<N$, put
  $\mathsf o_k=\sum_{r\le k}\mathsf g_r\mathsf E^-_{k-r}$ for $k<N$.
\item \emph{Left end.} Require $\mathrm{lo}(\mathsf X)>0$. Let $\varepsilon$ be an upper
  bound of $3e^{-50}/\mathrm{lo}(\mathsf X)$, and require $\varepsilon\le1$. With
  $L_1=\xia-\xiL$, put
  $P=e^{\ell L_1}+\mathrm{box}(2\varepsilon\,\mathrm{up}(e^{\ell L_1}))$ and $P'=\ell P$.
\item \emph{Middle steps.} For $j=0,1,\dots,n-1$ in turn, do the following, and then
  replace $(P,P')$ by the new values. For the point data and for the step data of
  step $j$, compute
  \[
  \begin{gathered}
  \mathsf D_0=\Lambda^{\rm b}+i\ell(\mathsf o_0+\mathcal G^{\rm b}\mathsf E^-_0),\qquad
  \mathsf D_k=i\ell(\mathsf o_k+\mathcal G^{\rm b}\mathsf E^-_k)\quad(1\le k<N),\\
  \mathsf R_0=1/\mathsf D_0,\qquad \mathsf R_k=-\mathsf R_0\sum_{r=1}^k\mathsf D_r\mathsf R_{k-r}\quad(1\le k<N),\\
  \mathsf V_k=i\ell\sum_{r\le k}\mathsf f_r\mathsf R_{k-r}+\ell^2[k=0]\quad(0\le k<N),
  \end{gathered}
  \]
  where $[k=0]$ is $1$ if $k=0$ and $0$ otherwise. Let $q_*=|\mathsf V^{\rm st}_0|^{\rm up}$,
  require that $q_*$ is finite, let $\kst$ be an upper bound of
  $(\max(q_*,\ell^2))^{1/2}$, require $\kst h_j\le700$, choose a dyadic number
  $M\ge\max(\kst|P|^{\rm up},|P'|^{\rm up})\,e^{\kst h_j}$, put
  $P^{\rm st}=P+\mathrm{box}(h_jM)$ and $P'^{\rm st}=P'+\mathrm{box}(h_jq_*M/\kst)$,
  and compute $\Psi_0=P$, $\Psi_1=P'$ and
  $\Psi_{k+2}=\frac1{(k+1)(k+2)}\sum_{r\le k}\mathsf V^{\rm pt}_r\Psi_{k-r}$, and in the same
  way $\Psi^{\rm st}_k$ from $P^{\rm st}$, $P'^{\rm st}$ and $\mathsf V^{\rm st}$, for
  $k\le N+1$. The new values are
  \[
  P_{\rm new}=\sum_{k<N}\Psi_kh_j^k+\Psi^{\rm st}_Nh_j^N,\qquad
  P'_{\rm new}=\sum_{k=1}^Nk\Psi_kh_j^{k-1}+(N+1)\Psi^{\rm st}_{N+1}h_j^N .
  \]
\item \emph{Right end and output.} With $L_2=\xiR-\xib$ and
  $W=\max(\ell|P|^{\rm up},|P'|^{\rm up})$, return
  \[
  \Evans^{\rm b}=e^{-\ell(\xiR-\xiL)}\Big(e^{\ell L_2}(P'+\ell P)
  +\mathrm{box}\big(4\varepsilon W\,\mathrm{up}(e^{\ell L_2})\big)\Big).
  \]
\end{enumerate}

\begin{proposition}
\label{prop:enclosure}
Let $\ell\ge1$ be an integer, let $\Lambda\subset\{\Re\lambda>0\}$ be a compact
axis-parallel box, and let $\Cell\subset[0,\infty)$ be a compact interval. Assume
that every ball operation satisfies \textup{(CP)}. Run the procedure of
Section~\ref{app:procedure} with $\ell$, a box $\Lambda^{\rm b}\supset\Lambda$ and a ball
$\mathcal G^{\rm b}\supset\Cell$. If it does not fail, it returns a box
$\Evans^{\rm b}$ of finite radius such that $\Evans_\ell(\lambda;G)\in\Evans^{\rm b}$ for every
$\lambda\in\Lambda$ and every $G\in\Cell$.
\end{proposition}
\begin{proof}
Fix $\lambda\in\Lambda$ and $G\in\Cell$, and let $\psi=\psi_\ell(\cdot;\lambda,G)$. We
show that every computed ball contains a corresponding exact quantity, defined
from $\lambda$, $G$ and $\chi$; the last one is $\Evans_\ell(\lambda;G)$. By (CP) it
suffices to check, at each step, that the exact quantity is the same function of
the earlier exact quantities, or else to justify the enclosure separately. For a
smooth function $F$ we write $F_k(\xi)=F^{(k)}(\xi)/k!$.

\emph{Step \textup{(i)}.} On $[\xia,\xib]$, $\chi=1$ and, $\chi$ being smooth, all its
derivatives vanish there, so $\prof^{(k)}=(AB)^{(k)}$, where
$A(\xi)=H(\frac{\xi-a_0}d)$, $B(\xi)=H(-\frac\xi d)$, $\bar A=1-A=H(\frac{a_0-\xi}d)$ and
$\bar B=1-B=H(\frac\xi d)$, as $1-H(x)=H(-x)$. Since $H'=2H(1-H)$, $A'=\frac2dA\bar A$, $\bar A'=-A'$,
$B'=-\frac2dB\bar B$ and $\bar B'=-B'$. The Leibniz rule gives
$A_{k+1}=\frac2{d(k+1)}\sum_{r\le k}A_r\bar A_{k-r}$, likewise for $B$, and
$\prof_k=\sum_{r\le k}A_rB_{k-r}$, $(\prof')_k=(k+1)\prof_{k+1}$,
$(e^{\pm2\xi})_k=(\pm2)^ke^{\pm2\xi}/k!$; the function $q=\prof e^{2\xi}$ has
$q_k=\sum_{r\le k}\prof_r(e^{2\xi})_{k-r}$. For $\xi\in S$, since $H$ is increasing,
$A(\xi)\in\mathsf A_0$, $\bar A(\xi)\in\bar{\mathsf A}_0$, $B(\xi)\in\mathsf B_0$ and
$\bar B(\xi)\in\bar{\mathsf B}_0$; by induction on $k$, each of
$\mathsf A_k,\bar{\mathsf A}_k,\mathsf B_k,\bar{\mathsf B}_k,\mathsf w_k,
\mathsf f_k,\mathsf E^\pm_k,\mathsf q_k$ contains the Taylor coefficient at $\xi$ of
$A,\bar A,B,\bar B,\prof,\prof',e^{\pm2\xi},q$ respectively. So the point data enclose
the Taylor coefficients at $s_j$, and the step data at every point of
$[s_j,s_{j+1}]$.

The function $g(\xi)=\int_{-\infty}^\xi\prof(\xi')e^{2\xi'}\dd\xi'$ has $g'=q\ge0$, so it is
nondecreasing. On $(-\infty,\xia]$, $0\le\prof\le A\le e^{2(\xi-a_0)/d}$, so
$0\le g(\xia)\le\int_{-\infty}^{\xia}e^{2(\xi-a_0)/d+2\xi}\dd\xi=e^{2x_a+2\xia}/(2/d+2)$,
that is, $g(s_0)\in\gamma_0$. Taylor's theorem with the Lagrange remainder,
applied to $t\mapsto\int_{s_j}^{s_j+t}q$, gives
\[
g(s_{j+1})-g(s_j)=\sum_{k<N}q_k(s_j)\frac{h_j^{k+1}}{k+1}+q_N(\eta_j)\frac{h_j^{N+1}}{N+1}
\]
for some $\eta_j\in[s_j,s_{j+1}]$. This increment is nonnegative and lies in the
ball in parentheses defining $\gamma_{j+1}$, so $g(s_j)\in\gamma_j$ for every $j$ by
induction, and $g(\xi)\in\mathrm{hull}(\gamma_j,\gamma_{j+1})$ for
$\xi\in[s_j,s_{j+1}]$ since $g$ is nondecreasing. As $g_r=q_{r-1}/r$ for $r\ge1$, the
balls $\mathsf o_k$ contain the Taylor coefficients of $ge^{-2\xi}$, at $s_j$ for the
point data and at every point of the step for the step data.

\emph{Step \textup{(ii)}.} Since $\lambda\in\Lambda^{\rm b}$,
$\Re\lambda\ge\mathrm{lo}(\mathsf X)>0$, so $\varepsilon\ge\epsend$ (as in
Lemma~\ref{lem:end-regions}) and, as $\varepsilon\le1$,
$e^{\epsend}-1\le e^\varepsilon-1\le(e-1)\varepsilon\le2\varepsilon$.
Lemma~\ref{lem:end-regions}(2) on $[\xiL,\xia]$, where $\psi_0(\xi)=e^{\ell(\xi-\xiL)}$
and $|Y(\xiL)|_w=\ell$, gives $|\psi(\xia)-e^{\ell L_1}|\le2\varepsilon e^{\ell L_1}$ and
$|\psi'(\xia)-\ell e^{\ell L_1}|\le2\ell\varepsilon e^{\ell L_1}$. Hence
$\psi(\xia)\in P$ and $\psi'(\xia)\in\ell P=P'$.

\emph{Step \textup{(iii)}.} Assume $\psi(s_j)\in P$ and $\psi'(s_j)\in P'$. On
$[\xia,\xib]$ the coefficient $V=\ell^2+i\ell\prof'/(\lambda+i\ell\Omega_G)$ is smooth, so
$\psi$ is smooth there. Put $D=\lambda+i\ell\Omega_G=\lambda+i\ell(g+G)e^{-2\xi}$ and
$R=1/D$, where $D\ne0$ as $\Re D=\Re\lambda>0$. The identity $DR=1$ gives
$R_0=1/D_0$ and $R_k=-R_0\sum_{r=1}^kD_rR_{k-r}$ for $k\ge1$, and
$V=\ell^2+i\ell\prof'R$ gives $V_k=i\ell\sum_{r\le k}(\prof')_rR_{k-r}+\ell^2[k=0]$. By
step (i) and (CP), $\mathsf V^{\rm pt}_k\ni V_k(s_j)$ and
$\mathsf V^{\rm st}_k\ni V_k(\xi)$ for every $\xi\in[s_j,s_{j+1}]$; in particular
$|V|\le q_*\le\kst^2$ on the step. For the a priori enclosure, put
$u(t)=\max(\kst|\psi(s_j+t)|,|\psi'(s_j+t)|)$ for $t\in[0,h_j]$. Then
$\kst|\psi(s_j+t)|\le\kst|\psi(s_j)|+\int_0^t\kst|\psi'|$ and
$|\psi'(s_j+t)|\le|\psi'(s_j)|+\int_0^t\kst\cdot\kst|\psi|$, so
$u(t)\le u(0)+\kst\int_0^tu$, and Gronwall's inequality gives
$u(t)\le u(0)e^{\kst t}\le M$. Hence, for $\xi\in[s_j,s_{j+1}]$,
$|\psi(\xi)-\psi(s_j)|\le h_jM$ and, since $|\psi|\le M/\kst$ on the step,
$|\psi'(\xi)-\psi'(s_j)|\le\int|V||\psi|\le h_jq_*M/\kst$; so $\psi(\xi)\in P^{\rm st}$
and $\psi'(\xi)\in P'^{\rm st}$. The equation $\psi''=V\psi$ gives
$(k+1)(k+2)\psi_{k+2}=\sum_{r\le k}V_r\psi_{k-r}$, so $\Psi_k\ni\psi_k(s_j)$ and
$\Psi^{\rm st}_k\ni\psi_k(\xi)$ for every $\xi\in[s_j,s_{j+1}]$ and $k\le N+1$.
Taylor's theorem with the integral remainder gives
\[
\psi(s_{j+1})=\sum_{k<N}\psi_k(s_j)h_j^k+h_j^N\int_0^1N(1-\tau)^{N-1}\psi_N(s_j+\tau h_j)\dd\tau .
\]
The weight $N(1-\tau)^{N-1}$ is a probability density on $[0,1]$, so the integral
lies in the closed convex hull of $\{\psi_N(\xi):\xi\in[s_j,s_{j+1}]\}$, which lies in
the closed convex box $\Psi^{\rm st}_N$. Hence $\psi(s_{j+1})\in P_{\rm new}$; the same
argument for $\psi'$, with $(\psi')_r=(r+1)\psi_{r+1}$, gives
$\psi'(s_{j+1})\in P'_{\rm new}$. By induction on $j$, $\psi(\xib)\in P$ and
$\psi'(\xib)\in P'$ at the end of step (iii).

\emph{Step \textup{(iv)}.} Apply Lemma~\ref{lem:end-regions}(2) on $[\xib,\xiR]$.
Since $(\psi_0'+\ell\psi_0)'=\ell(\psi_0'+\ell\psi_0)$,
$(\psi_0'+\ell\psi_0)(\xiR)=e^{\ell L_2}(\psi'+\ell\psi)(\xib)$, and
\[
|(\psi'+\ell\psi)(\xiR)-e^{\ell L_2}(\psi'+\ell\psi)(\xib)|\le2|Y(\xiR)-Y_0(\xiR)|_w
\le2e^{\ell L_2}|Y(\xib)|_w\cdot2\varepsilon\le4\varepsilon We^{\ell L_2}.
\]
Multiplying by $e^{-\ell(\xiR-\xiL)}$ gives $\Evans_\ell(\lambda;G)\in\Evans^{\rm b}$.
\end{proof}

\subsection{Certified Winding Numbers}

A set $B\subset\C$ is of \textbf{type P} if it does not meet $(-\infty,0]$, and
of \textbf{type N} if it does not meet $[0,\infty)$.

\begin{proposition}
\label{prop:winding}
Let $\Gamma$ be the boundary of a closed rectangle $\Rect'$ in $\{\Re\lambda>0\}$,
traversed counterclockwise and cut into segments $[p_i,p_{i+1}]$,
$i=0,\dots,I-1$, with $p_I=p_0$. Let $\Cell$ be a set, let $E(\cdot;G)$ be continuous
on $\Gamma$ for each $G\in\Cell$, and let $\mathcal B_i\subset\C$ be sets with
$E(\lambda;G)\in\mathcal B_i$ for all $\lambda\in[p_i,p_{i+1}]$ and $G\in\Cell$. Assume that
each $\mathcal B_i$ is of type P or of type N, and assign to each $i$ one type for which this
holds. Then $E(\cdot;G)\ne0$ on $\Gamma$ and, for every $G\in\Cell$,
\[
\wind\big(E(\Gamma;G),0\big)=\sum_{i\ \text{of type N}}\Big(\mathbf 1\big[\Im E(p_{i+1};G)<0\big]-\mathbf 1\big[\Im E(p_i;G)<0\big]\Big).
\]
In this sum the two terms at a point shared by two segments of type N cancel.
\end{proposition}
\begin{proof}
Fix $G\in\Cell$. Every $\mathcal B_i$ misses $0$, so $E(\cdot;G)\ne0$ on $\Gamma$. The
principal argument $\mathrm{Arg}\colon\C\setminus\{0\}\to(-\pi,\pi]$ is continuous on
$\C\setminus(-\infty,0]$, and
$\theta_{\rm N}(z)=\mathrm{Arg}(z)+2\pi\,\mathbf 1[\Im z<0]$ is continuous on
$\C\setminus[0,\infty)$, with values in $(0,2\pi)$. Since $E(\cdot;G)$ maps the
$i$-th segment into $\mathcal B_i$, $\mathrm{Arg}(E(\cdot;G))$ is a continuous
argument on a segment of type P, and $\theta_{\rm N}(E(\cdot;G))$ on one of type N. Each differs by a constant from a continuous argument along the
whole curve, so the winding number is $\frac1{2\pi}$ times the sum over the
segments of their increments. The increments of $\mathrm{Arg}$ telescope around the
closed curve to $0$, and what remains is the indicator sum. At a point shared by
two segments of type N the indicator enters once with each sign.
\end{proof}

\subsection{The Two Computations}
\label{app:runs}

\begin{definition}
\label{def:Nprime}
We consider the following two statements.
\begin{enumerate}[label=\textup{(N\arabic*$'$)}]
\item For every $G\in[0,\Gmax']$, $\Evans_3(\cdot;G)\ne0$ on $\partial\Rect$ and
  $\wind(\Evans_3(\partial\Rect;G),0)=1$.
\item There are $236$ closed intervals $\Cell_i$, $1\le i\le236$, with
  union $[0,\Gmax']$, and closed rectangles
  $\Rect_1,\dots,\Rect_{236}\subset\Rect$, the left edge of
  each of which lies on the line $\Re\lambda=\frac{211}{2500}$, such that for
  every $i$ and every $G\in\Cell_i$, $\Evans_3(\cdot;G)\ne0$ on $\partial\Rect_i$
  and $\wind(\Evans_3(\partial\Rect_i;G),0)=1$.
\end{enumerate}
\end{definition}

\begin{lemma}
\label{lem:Nprime-N}
\textup{(N1$'$)} implies \textup{(N1)}. \textup{(N1)} and \textup{(N2$'$)} imply
\textup{(N2)}.
\end{lemma}
\begin{proof}
Let $G\in[0,\Gmax]\subset[0,\Gmax']$. By
Lemma~I.\ref{I-lem:evans-basic}\textup{(1)} and its last assertion, $\Evans_3(\cdot;G)$ is holomorphic on
$\{\Re\lambda>0\}$, which contains $\Rect$ and every $\Rect_i$.

Assume \textup{(N1$'$)}. By the argument principle, $\Evans_3(\cdot;G)$ has exactly
one zero, counted with multiplicity, in the interior of $\Rect$. By
Lemma~\ref{lem:apriori}, every zero with $\Re\lambda\ge\frac9{200}$
lies in $\Rect$, and by \textup{(N1$'$)} none lies on $\partial\Rect$; so the zeros in
$\{\Re\lambda\ge\frac9{200}\}$ are exactly those in the interior of $\Rect$. This is
\textup{(N1)}.

Assume \textup{(N1)} and \textup{(N2$'$)}, let $\lamst$ be the zero of
\textup{(N1)}, and choose $i$ with $G\in\Cell_i$. By \textup{(N2$'$)} and the argument
principle, $\Evans_3(\cdot;G)$ has a zero $\lambda'$ in the interior of $\Rect_i$, whose left edge
lies on $\{\Re\lambda=\frac{211}{2500}\}$, so $\Re\lambda'>\frac{211}{2500}>\frac9{200}$, and $\lambda'=\lamst$ by \textup{(N1)}.
So the only zero of $\Evans_3(\cdot;G)$ in $\{\Re\lambda\ge\frac9{200}\}$ has real
part greater than $\frac{211}{2500}$, and there is none in
$\{\frac9{200}\le\Re\lambda\le\frac{211}{2500}\}$. This is \textup{(N2)}.
\end{proof}

Both computations use the procedure of Section~\ref{app:procedure} with $\ell=3$.
For a closed interval $\Cell$ with rational endpoints, a closed rectangle
$\Rect'=[x_0,x_1]\times[y_0,y_1]\subset\{\Re\lambda>0\}$ with rational vertices, a
number $\nu>0$ and integers $J,K\ge1$, the \textbf{certification} of
$(\Cell,\Rect')$ with density $\nu$, at most $J$ rounds and at most $K$ segments
is the following computation. The boundary of $\Rect'$ is traversed
counterclockwise from $(x_0,y_0)$, and each edge, of length $L$, is cut into
$\max(2,\lceil\nu L\rceil)$ segments of equal length; all cut points are rational.
For a segment $[p,q]$, the procedure is run with
$\Lambda^{\rm b}=(X_p\cup X_q)+i(Y_p\cup Y_q)$, where $X_p,X_q,Y_p,Y_q$ are the balls
that contain $\Re p$, $\Re q$, $\Im p$, $\Im q$, and with $\mathcal G^{\rm b}$ the
union of the balls that contain the two endpoints of $\Cell$; by (CP),
$\Lambda^{\rm b}\supset[p,q]$ and $\mathcal G^{\rm b}\supset\Cell$. The resulting box
$\Evans^{\rm b}$ is accepted as of type P if it is certified that
$\Im\Evans^{\rm b}>0$, or $\Im\Evans^{\rm b}<0$, or $\Re\Evans^{\rm b}>0$; otherwise it
is accepted as of type N if it is certified that $\Re\Evans^{\rm b}<0$. The first
round treats every segment; each later round treats the two halves, cut at the
midpoint, of every segment not accepted in the previous round, including those for
which the procedure failed. A segment is cut only if there are at most $K$
segments at that moment, and at most $J$ rounds are made; otherwise the
certification does not succeed. If every segment has been accepted, then at every
point $p$ where a segment of type N meets a segment of type P the procedure is run
with $\Lambda^{\rm b}$ the box of the balls that contain $\Re p$ and $\Im p$ and
with the same $\mathcal G^{\rm b}$, and the
indicator $\mathbf 1[\Im\Evans_3(p;G)<0]$ is taken to be $1$ if the resulting box
is certified to have negative imaginary part and $0$ if it is certified to have
positive imaginary part. The certification \textbf{succeeds} if every segment is
accepted and every such indicator is decided; its result is then the indicator sum
of Proposition~\ref{prop:winding}.

A \textbf{run} with at most $D$ halvings certifies a list of pairs
$(\Cell,\Rect')$: if the certification of $(\Cell,\Rect')$ does not succeed, $\Cell$
is halved at its midpoint and each half is certified with the same $\Rect'$, at
most $D$ times in succession; the certified pieces take the place of $\Cell$, and a
piece still not certified after $D$ halvings is reported as failed.

The run for \textup{(N1$'$)} uses the rectangle $\Rect$ and the $25$ intervals
\[
\Big[\frac{j}{20}d,\frac{j+1}{20}d\Big]\quad(0\le j\le14),\qquad
\Big[\frac{30+j}{40}d,\frac{31+j}{40}d\Big]\quad(0\le j\le9),
\]
whose union is $[0,\Gmax']$, with density $\nu=20$, at most $16$ rounds, at most
$6000$ segments and at most $4$ halvings. The run for \textup{(N2$'$)} uses $236$
intervals $\Cell_1,\dots,\Cell_{236}$, numbered from left to right, and rectangles
$\Rect_i=[\frac{211}{2500},x_i]\times[y_i,y_i']$, with density
$\nu=3/(x_i-\frac{211}{2500})$, at most $10$ rounds, at most $800$ segments and at
most $5$ halvings. These $236$ pairs, with every number an exact fraction, form the file
\texttt{sector3\_gap\_cells.json} of the reproduction package
(Section~\ref{app:reproduce}), and the file \texttt{sector3\_count\_cells.json}
lists the $25$ intervals of the first run, also numbered from left to right. In each of these lists the endpoints of
the intervals are rational, the first interval starts at $0$, the last ends at
$d=\Gmax'$, and consecutive intervals share an endpoint; every $\Rect_i$ has
rational vertices, satisfies $x_i>\frac{211}{2500}$ and $y_i<y_i'$, and lies in
$\Rect$. The re-check of Section~\ref{app:reproduce} confirms these facts, and that
every rectangle of the first run is $\Rect$, for the output files of the runs, in
exact rational arithmetic.

\begin{proposition}
\label{prop:runs}
\textup{(N1$'$)} and \textup{(N2$'$)} hold.
\end{proposition}
\begin{proof}
The proof is the computation just described, carried out under the trusted base of
Section~\ref{app:arith}. Each of the $25$ certifications of the first run and each
of the $236$ of the second succeeded at once, with result $1$, so no interval was
halved and the output files list exactly the intervals and rectangles described
above.

Consider a certification of $(\Cell,\Rect')$ that succeeds with result $1$. By
Lemma~I.\ref{I-lem:evans-basic}\textup{(1)} and its last assertion, $\Evans_3(\cdot;G)$ is continuous on
$\{\Re\lambda>0\}\supset\partial\Rect'$ for every $G\in\Cell$. By
Proposition~\ref{prop:enclosure}, the box of each accepted segment $[p,q]$
contains $\Evans_3(\lambda;G)$ for all $\lambda\in[p,q]$ and $G\in\Cell$, and each decided
indicator equals $\mathbf 1[\Im\Evans_3(p;G)<0]$ for every $G\in\Cell$; by (CP) for
the comparisons, a box accepted as of type P misses $(-\infty,0]$ and one accepted
as of type N misses $[0,\infty)$. The accepted segments are consecutive and cover
$\partial\Rect'$. By Proposition~\ref{prop:winding}, $\Evans_3(\cdot;G)\ne0$ on
$\partial\Rect'$ and $\wind(\Evans_3(\partial\Rect';G),0)=1$ for every $G\in\Cell$.
Applied to the first run, whose intervals cover $[0,\Gmax']$, this gives
\textup{(N1$'$)}; applied to the second, with the properties of the $\Cell_i$ and
$\Rect_i$ stated above, it gives \textup{(N2$'$)}.
\end{proof}

\subsection{The Closed-Form Inequality}

\begin{proposition}
\label{prop:N3}
\textup{(N3)} holds.
\end{proposition}
\begin{proof}
We first bound $\epsst$. Since $H(x)\le e^{2x}$, Equation~\eqref{eq:xayb} gives
$\epsst\le e^{2\max(x_a,y_b)}<e^{-49}<10^{-21}$, as $e^{49}>2.7^{49}>10^{21}$
(proof of Lemma~\ref{lem:apriori}).

We next bound $e^{a_0}$. From Definition~I.\ref{I-def:ring}, $a_0=-7304210039150669\cdot2^{-54}<-0.4054$. The partial
sum $\sum_{k\le5}0.4054^k/k!$ of the exponential series is a lower bound for
$e^{0.4054}$, and it exceeds $1.49989$, so $e^{a_0}<e^{-0.4054}<1/1.49989<0.66672$.
Squaring and multiplying, $0.66672^2<0.444516$, $0.66672^3<0.296368$,
$0.66672^6<0.087834$ and $0.66672^9<0.026032$.

We now bound $\Thl{\ell}$ for $\ell\in\{6,9\}$. Put $u=\pi\ell d$. Since $\pi<3.1416$ and $0<d<0.0200001$, we have $0<u<0.377$ for $\ell=6$ and
$0<u<0.5655$ for $\ell=9$; in particular $0<u<\pi$. For $u\ge0$,
$\sin u\ge u-u^3/6$, since the difference vanishes at $0$ and has derivative
$\cos u-1+u^2/2$, which is nonnegative (it vanishes at $0$ and its derivative $u-\sin u$ is
nonnegative). So $\Thl{\ell}=u/\sin u\le(1-u^2/6)^{-1}$ as
$u^2<6$. Hence $\Thl{6}\le(1-0.377^2/6)^{-1}<(1-0.023689)^{-1}<1.024264$
and $\Thl{9}<(1-0.0533)^{-1}<1.0564$. Therefore
\[
\begin{gathered}
e^{6a_0}\Thl{6}+\epsst<0.087834\cdot1.024264+10^{-21}<0.08997,\\
e^{9a_0}\Thl{9}+\epsst<0.026032\cdot1.0564+10^{-21}<0.0276,
\end{gathered}
\]
and both are less than $\frac9{100}$, which is \textup{(N3)}.
\end{proof}

\begin{proof}[Proof of Theorem~\ref{thm:N-holds}]
By Proposition~\ref{prop:runs}, \textup{(N1$'$)} and \textup{(N2$'$)} of
Definition~\ref{def:Nprime} hold; by Lemma~\ref{lem:Nprime-N}, \textup{(N1)} holds, and
then \textup{(N2)}; by Proposition~\ref{prop:N3}, \textup{(N3)} holds.
\end{proof}

\subsection{Reproduction}
\label{app:reproduce}

The computations of Section~\ref{app:runs} can be repeated with version~1.0.0 of
the package \texttt{euler-ring-certificate}, which is available under the Apache
License 2.0 in the public repository \url{https://github.com/mcoiculescu/NonUniqueEuler}
and contains the code, the two
certificate files of Section~\ref{app:runs} and the reference outputs. It requires
only python-flint~0.9.0, which bundles FLINT~3.6.0, under CPython~3.10 or later.
The program refuses to run with assertions disabled, and, unless explicitly
allowed, with another FLINT version. From the top directory of the package, the
command \mbox{\texttt{./reproduce.sh 8}} runs, with $8$ worker processes, among other
items: an environment check; the unit tests; the certification of the $25$
intervals of \textup{(N1$'$)} and of the $236$ intervals and rectangles of
\textup{(N2$'$)} (commands \texttt{ringcert sector3-count} and
\texttt{ringcert sector3-gap}); a re-check of the outputs in exact rational
arithmetic (\texttt{ringcert check}); and a comparison with the reference
outputs. A successful run reports all items as passed.

The outputs record, for every certified interval, its endpoints, its rectangle,
its result, and the numbers of segments and indicators. From the output files
alone, the re-check confirms that the intervals of each run cover $[0,\Gmax']$,
with consecutive intervals sharing an endpoint, that every result is $1$, that no
indicator is undecided and no interval failed, that every rectangle of the first
run is $\Rect$, and that every rectangle of the second run is nondegenerate, has
left edge on the line $\Re\lambda=\frac{211}{2500}$ and lies in $\Rect$; it also
checks that these output files are the ones recorded in the summary of the run.

The output files \texttt{sector3\_count.json} and \texttt{sector3\_gap.json}
contain no timing or environment data; on an Apple arm64 machine with
CPython~3.13.5 they are byte-identical to the reference outputs of the package.
The SHA-256 digests of the certificate files and of these outputs are the following.
\begin{center}
\scriptsize
\begin{tabular}{ll}
\texttt{sector3\_count\_cells.json} &
\texttt{518a0ca4f05fd24bae273aebea20862fad1c1beb46dbdee9f8b5fb76c8729dc1}\\
\texttt{sector3\_gap\_cells.json} &
\texttt{fb95a3f72eee1f017ddc5c26e7dfca52841759603e34b0131a440c8992086c91}\\
\texttt{sector3\_count.json} &
\texttt{2aa43b18d7289b7288f4bc660b1fdbec547458949ee675ef4b23f65cbcf1a1c9}\\
\texttt{sector3\_gap.json} &
\texttt{f06feec05cc6633b390bb9ee6191ce59c3bd5231a01ffa0f48263c8c9173a786}
\end{tabular}
\end{center}
The reference outputs were produced on that machine only. With another FLINT
version or on other hardware the numbers of segments may differ, and an interval
may be halved. The proof needs only that the re-check passes, that is, that the output files
have the properties listed above. If some intervals
were halved, the argument of Proposition~\ref{prop:runs} gives \textup{(N1$'$)},
and \textup{(N2$'$)} with the certified pieces, each paired with the rectangle of
its interval, in place of the $236$ intervals; the proof of
Lemma~\ref{lem:Nprime-N} does not use their number.
